\chapter{<<vi:Khái niệm cơ bản||en:Basic concepts>>}
\label{ch:basics}

%%IF box
\begin{summary}[<<vi:Trong chương này||en:In this chapter>>]
\begin{itemize}
  \item <<vi:Ý thứ nhất||en:First idea>>
  \item <<vi:Ý thứ hai||en:Second idea>>
\end{itemize}
\end{summary}
%%ENDIF

\section{<<vi:Định nghĩa||en:Definitions>>}
<<vi:Từ cần đưa vào chỉ mục được đánh dấu bằng||en:Terms for the index are marked with>> \verb|\index{...}|<<vi:, ví dụ hàm lồi||en:, for example convex function>>\index{<<vi:hàm lồi||en:convex function>>}.
%%IF math
\begin{definition}
<<vi:Một hàm||en:A function>> $f:\R^n\to\R$ <<vi:là lồi nếu||en:is convex if>>
$f(\lambda x + (1-\lambda) y) \le \lambda f(x) + (1-\lambda) f(y)$.
\end{definition}

\begin{theorem}
<<vi:Mọi cực tiểu địa phương của hàm lồi đều là cực tiểu toàn cục.||en:Every local minimum of a convex function is a global minimum.>>
\end{theorem}
\begin{proof}
\todo{<<vi:chứng minh||en:proof>>}
\end{proof}
%%ENDIF

%%IF box
\begin{examplebox}
<<vi:Ví dụ minh họa cho định nghĩa.||en:An example of the definition.>>
\end{examplebox}

\begin{tip}
<<vi:Mẹo nhỏ cho người đọc.||en:A small tip for the reader.>>
\end{tip}
%%ENDIF

\section{<<vi:Bài tập||en:Exercises>>}
\begin{enumerate}
  \item \todo{<<vi:bài tập 1||en:exercise 1>>}
\end{enumerate}
